\subsection{Two-saddle comparison and action matching}
\label{sec:two-saddle-comparison}

On this coefficient face the two incident saddles lie on different
scales, and both characteristic ratios tend to zero. Their local maps
must be compared after the intervening nonlinear flight. We first derive
its scalar action equation and the actual local section coordinates.
Section~\ref{subsec:deg-degree-drop} then constructs their physical
receiver and follows the passage over the full action range.

\subsection{The two saddle scales}

Let \(0<\kappa,\tau\ll1\), \(\delta=\kappa\tau\). In coordinates
which put the negative saddle at \(u=-1\), the other two equatorial
roots are \(u=a\) and \(u=c\). Write
\begin{equation}\label{fx:pd}
 p_{\rm d}=-\mathcal R(u+1)(u-a)(u-c),
\end{equation}
where
\begin{equation}\label{fx:coefficient-graph}
 a=\kappa A_{\rm d},\qquad a+c=1-3\nu .
\end{equation}
The quantities \(A_{\rm d},\mathcal R,c,c-a,1-3\nu\) and their
reciprocals belong to fixed positive compact sets; \(A,C\) vary in
fixed compact sets. Thus \(a\) tends to zero while \(-1\) and \(c\)
remain separated. The vector field has the form
\[
 \dot u=\delta p_{\rm d}-Kw,\qquad
 \dot w=-w(L+Ew),
\]
with
\begin{equation}\label{fx:coefficients}
\begin{aligned}
 K&=1-\delta\mathcal R A(u+\nu)
              +C\delta^2\mathcal R^2(u+\nu)^2,\\
 L&=u\{1+\delta\mathcal R(u+2\nu)\},\qquad
 E=C\delta^2\mathcal R^2(u+\nu).
\end{aligned}
\end{equation}
On the intervening axis segment put \(h=-\delta p_{\rm d}>0\).
Positive time runs from the cut \(u_+=a(1-\eta_+)\) to
\(u_-=-1+\eta_-\), with fixed \(0<\eta_\pm<1\). The orbit equation is
\begin{equation}\label{fx:graph}
 w_u=\frac{w\{L(u)+E(u)w\}}{h(u)+K(u)w}.
\end{equation}

The tangential eigenvalues at \(a\) and \(-1\) are
\(\delta\ell_+\) and \(-\delta\ell_-\), whereas their normal
eigenvalues are \(-b_+\) and \(b_-\), where
\[
\begin{aligned}
 \ell_+&=\mathcal R(1+a)(c-a),&
 b_+&=a\{1+\delta\mathcal R(a+2\nu)\},\\
 \ell_-&=\mathcal R(1+a)(1+c),&
 b_-&=1-\delta\mathcal R(1-2\nu).
\end{aligned}
\]
Define the small characteristic ratios used in the local maps by
\begin{equation}\label{fx:exponents}
 \chi=\delta\ell_+/b_+,\qquad
 \beta=\delta\ell_-/b_-,\qquad q=\beta/\chi .
\end{equation}
Thus \(\chi\) is the exponent for the inverse positive Dulac map,
whereas \(\beta\) is the exponent for the forward negative Dulac map.
The normal rate at the positive saddle is of order \(\kappa\),
while that at the negative saddle is bounded away from zero. Hence
\(\chi\) is of order \(\tau\), \(\beta\) is of order
\(\kappa\tau\), and \(q=\kappa Q_{\rm d}\) with
\(Q_{\rm d}\) a positive unit. In particular,
\begin{equation}\label{fx:log-smallness}
 \beta|\log a|+q|\log a|\longrightarrow0.
\end{equation}
The inequality \(\beta\ll\chi\) will determine the sign in the
comparison of the two Dulac maps. To transfer that sign to the return,
we also need the nonlinear transition between their local sections.

\subsection{A transition in accumulated action}

Writing \(A_m=L/h\) separates the linear normal equation in
\eqref{fx:graph} from its nonlinear part. The following calculation
also applies to general coefficients for which \(h>0\):
\begin{equation}\label{eq:comparison-middle}
 w_u=A_m(u)w+\frac{E(u)-A_m(u)K(u)}{h(u)+K(u)w}\,w^2,\qquad h>0 .
\end{equation}
Let \(J'=A_mJ,\ J(u_+)=1\), and write \(h_-=h(u_-),J_-=J(u_-)\).
For a fixed \(\lambda>0\), put
\[
 \rho=\frac{h_-J}{J_-h},\qquad
 w=\lambda h\rho V .
\]
The factor \(J/J_-\) removes the linear amplification between the
cuts, and \(\rho(u_-)=1\). This normalization isolates the remaining
nonlinear change of height. In the saddle configuration above the
amplification is exponentially large in \(1/\beta\); its exact value
is given by the partial fractions in \eqref{fx:partial-fractions}.
Direct substitution retains the denominator in the equation for \(V\):
\begin{equation}\label{eq:comparison-action}
 dV=-\frac{\lambda KV^2}{1+\lambda K\rho V}\,d\rho
 +\frac{\lambda\rho(E-Kh'/h)V^2}{1+\lambda K\rho V}\,du .
\end{equation}

For bounded \(V=V(u_+)\), an early negative cut gives the actual
section coordinates before the nonlinear amplification layer. Larger
values of \(V\) require a different description of the same passage.
Section~\ref{sec:fixed-exit-comparison}
uses the coordinate \(\Theta=\beta\log V\) on a bounded nonnegative
interval. For a positive limiting \(\Theta\) in that interval,
\(V\) is exponentially large in \(1/\beta\), and the orbit leaves
the small-height region before the negative saddle. Its height on a
fixed exit section then supplies a regular coordinate. These are two
overlapping parts of the orbit interval for the same coefficients.

\subsection{Actual section coordinates at both stopped endpoints}
\label{subsec:actual-dulac-matching}

For bounded amplified height, the middle passage is close to its linear
normalization up to an early negative cut. We stop there and use exact
saddle invariants to continue to both actual sections. This retains the
nonlinear section maps needed to compare the two small powers.

\begin{lemma}[Uniform actual-section matching]
\label{lem:actual-dulac-matching}
Consider \eqref{fx:pd}--\eqref{fx:coefficients}, with the compact
coefficient conditions stated there. Fix \(V_1<\infty\) and let
\(V\) be the raw middle input in \eqref{eq:comparison-action}.
Choose the local cuts sufficiently small in the order in the proof,
before decreasing \(\kappa,\tau\).
On \(0\le V\le V_1+1\) there is an increasing change
\(\widehat V=V+O_{C^3}(\epsilon_\theta)\), fixing zero, where
\(\epsilon_\theta/\chi\to0\). The two actual section coordinates,
including the finite closing map on the negative side, have the form
\begin{equation}\label{eq:actual-dulac-pair}
 x_{\rm in}=x_b+P_*\Psi(\widehat V^\chi),\qquad
 f_\theta(mt_{\rm out})=b_*+a_*D
       \bigl(\beta^{-\beta}C(\widehat V)^\beta\bigr),
 \qquad P_*,a_*>0,
\end{equation}
for \(m>0\). Put \(e_C(v)=vC'(v)/C(v)>0\) for \(v>0\).
Uniformly for \(0<V\le V_1+1\),
\begin{equation}\label{eq:actual-dulac-bounds}
 \begin{gathered}
 C(0)=0,\quad C'(0)=1,\quad
 \|C-\mathrm{id}\|_{C^3}\le\epsilon_\theta,\qquad
 \left|\frac{V e_C'(V)}{e_C(V)}\right|=o(\chi),\\
 D'>c_D>0,\qquad |D''/D'|\le C_D,\qquad
 \sup_{0\le t\le(V_1+2)^\chi}
       |t\Psi''(t)/\Psi'(t)|\longrightarrow0,\qquad P_*\to0 .
 \end{gathered}
\end{equation}
The maps \(\widehat V,C,D,\Psi\) extend in \(C^3\) to their
zero endpoints; this does not assert ordinary \(C^3\) regularity
of the compositions with the small powers of \(\widehat V\).
The normalized estimates have a uniform one-sided extension as
\(j=1/c\downarrow0\) at the remote-root face. Actual raw section
maps and their positive affine scales are asserted only for \(j>0\).
\end{lemma}

\begin{proof}
We first construct local coordinates in which the saddle invariant is
exact and the derivative bounds remain uniform as its exponent vanishes.
After dividing by the positive normal rate, and reversing time at
the positive saddle, both fields on a fixed small box take the form
\begin{equation}\label{eq:saddle-normal-time}
 \dot x=-\rho x p(x)+c(x)y,\qquad
 \dot y=y\{1+k(x)y\},\qquad \rho\in\{\beta,\chi\}.
\end{equation}
Here \(p(0)=1\), \(p>c_p>0\), and \(p,c,k\) have uniform
bounds through any prescribed finite order. At the negative saddle
\(x=u+1,y=w\), \(c=K/L<0\); at the positive saddle
\(x=(u-a)/a,y=w/a^2\), \(c>0\). Both \(|c|\) and its
reciprocal are bounded. These assertions follow by substituting the
displayed scalings in \eqref{fx:coefficients}; in particular the
positive normal rate is divided out only after scaling by \(a\).

Put \(h=-\rho xp(x)\). Construct the footpoint \(\pi\) by backward
axis transport as in Lemma~\ref{lem:normal-time-projection}, and put
\(J=\int_0^\infty k(X(-T))Y(-T)\,dT\), using bounded extensions.
Here \(Y(-T)\le ye^{-T/2}\), \(h=O_{C^4}(\rho)\), and the
couplings \(c_xY,k_xY^2\) have integrable mass. The same triangular
argument bounds tangential rows by \(C(1+T)^Ne^{C\rho T}\) and
normal rows by \(C(1+T)^Ne^{-(1/2-C\rho)T}\).
For a normalized coefficient parameter \(\vartheta\), its new forcing
is \(-\bigl(h_\vartheta+c_\vartheta Y,k_\vartheta Y^2\bigr)\):
the tangential term adds polynomial factors, while normal terms keep a
decaying height factor. Thus the backward integrals and their mixed rows
through five phase derivatives and one parameter derivative have an
integrable majorant \(C(1+T)^Ne^{-(1/2-C\rho)T}\), after fixing the
finite order and then decreasing \(\rho\). These mixed rows use axis-flow
derivatives through order seven, supplied by the prescribed finite
coefficient bounds.

Consequently \(\dot\pi=h(\pi)\), and \(Z=ye^{-J}\) satisfies
\(\dot Z=Z\). The axis coordinate
\[
 \ell(x)=x\exp\int_0^x\frac{p(s)^{-1}-1}{s}\,ds
\]
has uniform finite-order bounds and positive derivative. Hence
\(\Xi=\ell(\pi)\) gives the exact invariant
\begin{equation}\label{eq:exact-saddle-invariant}
 \dot\Xi=-\rho\Xi,\qquad I=\Xi Z^\rho,\qquad \dot I=0.
\end{equation}
This is a uniformly regular local coordinate change. Since
\(\pi_y(x,y)=-c(x)+o_\rho(1)+O(y)\) and \(|c|\) is bounded below,
choose the height strip first and then \(\rho\) to preserve its sign.
All these bounds concern intrinsic coordinates, before extracting raw
section scales.

Fix a permissible small upper height \(y_{\max}\) in the negative
box. Choose \(x_->0\) small compared with this height, and then
\(\eta_0>0\) sufficiently small compared with both, as quantified
in the output construction below. Finally fix
\(0<\eta_-<\eta_0\) and put \(u_0=-1+\eta_0\).
This cut lies in the negative box, and
strictly before the cut \(u_-\) in positive time. If
\(J_m(u)=\exp\int_{u_+}^u L/h\), put
\[
 \varepsilon_+=\frac{\lambda h_-}{a^2J_m(u_-)},\qquad
 \varepsilon_0=\frac{\lambda h_-J_m(u_0)}{J_m(u_-)}.
\]
The raw heights at \(u_+\) and \(u_0\) are respectively
\(a^2\varepsilon_+V\) and \(\varepsilon_0 M_0(V)\).
Partial fractions \eqref{fx:partial-fractions} imply, for some
fixed \(c_0>0,N<\infty\),
\begin{equation}\label{eq:early-saddle-smallness}
 \varepsilon_++\varepsilon_0
       \le C\beta^{-N}e^{-c_0/\beta},\qquad
 c\le\varepsilon_0^\beta\le C,\qquad
 \varepsilon_+^\chi\le Ce^{-c_0/(2q)} .
\end{equation}
To check uniformity, \(\beta\le Ca\), so every inverse power of
\(a\) is bounded by a power of \(\beta^{-1}\). The logarithmic
terms in the partial fractions are
\(q\log a\), \(\beta\log a\), and \(\beta\log\beta\),
all tending to zero. The action between \(u_0\) and \(u_-\)
has a fixed positive limit, while the full action to \(u_-\)
has the positive limit in \eqref{fx:gamma-endpoint}.
These facts give all three estimates, with the last obtained by
multiplying the logarithm of the first by \(\chi=\beta/q\).

Before \(u_0\), the middle map satisfies
\begin{equation}\label{eq:early-middle-identity}
 M_0(0)=0,\quad M_0'(0)=1,\qquad
 \|M_0-\mathrm{id}\|_{C^3([0,V_1+1])}
       \le C\beta^{-N}e^{-c_0/(2\beta)}.
\end{equation}
Indeed \(\rho'/\rho=(L-h')/h\) has at most one zero, a minimum;
the endpoint actions give
\(\sup_{[u_0,u_+]}\rho\le C\beta^{-N}e^{-c_0/\beta}\).
The ordinary equation is
\[
 V_u=\mathscr F=\frac{\lambda\rho(E-KL/h)V^2}{1+\lambda K\rho V}.
\]
On \([0,V_1+2]\) its denominator is at least one, and the integral
of its \(C^3_V\) norm is at most
\(C\beta^{-N}(1+|\log\beta|)e^{-c_0/\beta}\).
The first-exit and triangular-variation argument of
Lemma~\ref{lem:scalar-normalization} proves the estimate;
\(\mathscr F=O(V^2)\) gives \(M_0'(0)=1\). No ordinary
\(u\)-derivative at the shrinking positive cut is needed.

On the positive middle cut write the coordinate restrictions as
\(\Xi_{+,m}(y),Z_{+,m}(y)\), and let
\(z_+=Z_{+,m}'(0)>0\). Define
\[
 \widehat V=
 \frac{Z_{+,m}(\varepsilon_+ V)}{z_+\varepsilon_+}
 \left(\frac{-\Xi_{+,m}(\varepsilon_+ V)}
                  {-\Xi_{+,m}(0)}\right)^{1/\chi} .
\]
The weak coordinate is negative on this cut. Taylor's formula and
\eqref{eq:early-saddle-smallness} give
\(\widehat V=V+O_{C^3}(\varepsilon_+/\chi)\), with
\(\widehat V'(0)=1\). Its inverse has the same bounds.
Choose a fixed incoming local cut \(x=x_+>0\) in the positive
box. The stopped orbit \(\Xi=0\) intersects it at a height
\(y_+>c>0\), since \(\Xi_y<0\). On a fixed chord above that
height the function
\[
 L_+(y)=-\Xi_{+,\rm out}(y)Z_{+,\rm out}(y)^\chi
\]
is increasing, vanishes at \(y_+\), and has a uniformly bounded
\(C^3\) inverse. The exact invariant gives
\[
 L_+(y)=\Sigma_+\widehat V^\chi,\qquad
 \Sigma_+=-\Xi_{+,m}(0)(z_+\varepsilon_+)^\chi
                    \le Ce^{-c_0/(2q)}.
\]

Continue this chord to the actual positive section. In
\((X,Y)=(u/a,w/a^2)\) the limiting stopped graph is
\(Y=(X^2-1)/2\), \(X\ge1+x_+\). A smaller chord has positive
height and denominator margins to a fixed large cut and then enters a
common parabolic band. Apply the outgoing-tail normalization
\eqref{ann:tail-first-derivative}--\eqref{eq:N-bound} and the
fixed-port admission argument of Lemma~\ref{ann:p0-fixed-tail-ports}.
Both use only this half-family and its common band. If \(d_+\) is the
exact raw derivative at \(y_+\), then
\(p=(x_{\rm in}-x_b)/d_+\) gives
\(L_+(y)=\Phi_+(p)\), \(\Phi_+(0)=0\), with positive derivative
units and uniform \(C^3\) inverse bounds. Retain the absolute gain
\(d_+=\Omega a^{2-\lambda_\gamma}U_+\le C\), where \(U_+\) is
a positive unit.
Set
\[
 \sigma_+=\Phi_+^{-1}(\Sigma_+),\qquad
 \Psi(t)=\frac{\Phi_+^{-1}(\Sigma_+t)}{\sigma_+},\qquad
 P_*=d_+\sigma_+ .
\]
Then \(\sup|t\Psi''/\Psi'|\le C\Sigma_+\), giving the
positive assertions without any bound on unscaled raw derivatives.

At the negative cut \(u_0\), denote the coordinate restrictions by
\(\Xi_{-,m}(w),Z_{-,m}(w)\), and put
\(\xi_0=\Xi_{-,m}(0)>0\), \(z_0=Z_{-,m}'(0)>0\).
For \(v\) in a fixed bounded interval set
\[
 \mathcal G(v)=\frac{Z_{-,m}(\varepsilon_0v)}{z_0\varepsilon_0}
 \left(\frac{\Xi_{-,m}(\varepsilon_0v)}{\xi_0}\right)^{1/\beta},
 \qquad C=\mathcal G\circ M_0\circ(V\mapsto\widehat V)^{-1}.
\]
This is the actual strong-coordinate correction on the early cut.
Its \(C^3\) difference from the identity is
\(O(\varepsilon_0/\beta)\); the cut \(\eta_0\) is fixed.
Together with \eqref{eq:early-middle-identity} this yields
\(C'(0)=1\) and the claimed estimate for \(C-\mathrm{id}\),
with \(\epsilon_\theta=C\beta^{-N}e^{-c_0/(3\beta)}\).
Hadamard division gives \(C(v)=v c(v)\),
\(\|c-1\|_{C^2}\le C\epsilon_\theta\), so
\(|v e_C'/e_C|\le C\epsilon_\theta=o(\chi)\).

On the fixed outer negative cut \(u=-1-x_-\),
\(\Xi_w>0\), so the stopped branch \(\Xi=0\)
meets it at a height \(w_->c x_->0\). On a slightly larger
positive-height chord its invariant
\[
 I_-(w)=\Xi_{-,\rm out}(w)Z_{-,\rm out}(w)^\beta
\]
has \(I_-'>c>0\) and uniformly bounded \(C^3\) norm.
Writing \(r=\beta^{-\beta}C(\widehat V)^\beta\), equality
of the two invariants is precisely
\[
 I_-(w)=\alpha_\theta r,\qquad
 \alpha_\theta=\xi_0(z_0\varepsilon_0\beta)^\beta .
\]
The constant \(\alpha_\theta\) is bounded above and away from
zero by \eqref{eq:early-saddle-smallness}. Also
\(\alpha_\theta\le C\eta_0\), while \(r\le1+o(1)\)
on the fixed amplified-height interval. The choices above can be
made with \(\pi(-x_-,y_{\max})\ge c y_{\max}>0\), and
then with
\[
 (1+\epsilon)\sup_\theta\alpha_\theta
       <\tfrac12\inf_\theta I_-(y_{\max}).
\]
The whole invariant interval consequently lies strictly inside the
output chord, with a uniform buffer. The inverse function theorem gives
\(w=D_0(r)\) on \([0,1+\epsilon]\), with uniform \(C^3\)
and inverse bounds. At \(r=0\) its output is \(w_->0\).
The extension at the excluded endpoint \(\widehat V=0\) is thus
obtained on a positive-height chord.

From this cut the limiting stopped graph is \(w=(u^2-1)/2\).
Its buffered chord has the same finite-cut margins and negative-tail
band. The preceding normalization and admission argument, now on the
negative half-tail, gives
\(T(w)=t_b+\Omega U_-B_-(w)\), with \(B_-'>c\) and uniform
\(C^3\) map and inverse bounds. In particular \(|T''/T'|\le C\),
although its raw slope may vanish.
For the divided finite closing map, the chain rule gives
\[
 \frac{(f_\theta(mT)/m)''}{(f_\theta(mT)/m)'}
   =\frac{T''}{T'}+
      m\frac{f_\theta''(mT)}{f_\theta'(mT)}T' .
\]
The last term is bounded on the fixed output window since
\(0<\Omega\le C\). Normalizing its exact value and positive
slope at a reference point, and composing with \(D_0\), gives
\(D\) and the constants \(b_*,a_*\) in
\eqref{eq:actual-dulac-pair}. This proves the negative assertions.

Finally, in the remote-root chart the coefficients on both local
boxes are regular functions of \(j\). The source-tail scales are
exactly \(\Omega\) and \(\Omega a^{2-\lambda_\gamma}\);
\(|\lambda_\gamma|<1\) keeps both bounded above. No reciprocal
bound on these raw scales was used. The fixed-cut action has no
\(\beta\log j\) term, so every estimate above is uniform as
\(j\downarrow0\), as a one-sided family of physical passages.
\end{proof}

